\begin{tabularx}{\textwidth}{@{}>{\raggedright\arraybackslash}p{0.20\textwidth} X r r@{}}
\toprule
Feature & Description & Win rate (95\% CI) & $q$ \\
\midrule
\bfseries Concrete illustration$^{*}$ & The argument makes its case through specific, vivid, or worked-through examples rather than abstraction alone. & 0.604 {\tiny$[0.570, 0.638]$} & $<$0.001 \\
\bfseries Qualified rebuttal$^{*}$ & The challenger narrows the claim, marks exceptions, or emphasizes contextual dependence instead of arguing absolutely. & 0.596 {\tiny$[0.562, 0.630]$} & $<$0.001 \\
\bfseries Structured rebuttal development$^{*}$ & The challenger builds an organized, multi-step, or multi-pronged case rather than a single bare point. & 0.583 {\tiny$[0.548, 0.617]$} & $<$0.001 \\
\bfseries Scope broadening & The exchange expands or contracts the discussion by moving between narrow cases and wider principles or issue areas. & 0.583 {\tiny$[0.548, 0.617]$} & $<$0.001 \\
\bfseries Counterexample use$^{*}$ & The challenger undermines the OP with exceptions, alternative cases, or edge-case scenarios that break the generalization. & 0.581 {\tiny$[0.547, 0.616]$} & $<$0.001 \\
\bfseries Challenger evidentiary support$^{*}$ & The challenger backs the rebuttal with external evidence, citations, factual corrections, or expert authority. & 0.575 {\tiny$[0.541, 0.609]$} & $<$0.001 \\
\bfseries Conceptual clarification$^{*}$ & The challenger disputes definitions, categories, meanings, or criteria in order to reframe the disagreement. & 0.574 {\tiny$[0.540, 0.609]$} & $<$0.001 \\
\bfseries Structural explanation & The challenger explains the issue through broader social, cultural, or systemic forces rather than isolated choices. & 0.574 {\tiny$[0.540, 0.609]$} & $<$0.001 \\
\bfseries Abstract or technical framing$^{*}$ & The exchange reasons at a more theoretical, philosophical, scientific, or specialized level rather than in everyday concrete terms. & 0.574 {\tiny$[0.539, 0.608]$} & $<$0.001 \\
\bfseries Analogical reasoning$^{*}$ & The argument persuades by comparing the case to parallel real or hypothetical cases, including extended analogies and thought experiments. & 0.570 {\tiny$[0.536, 0.604]$} & $<$0.001 \\
\bfseries Mechanistic explanation$^{*}$ & The challenger explains or critiques the causal mechanism by which the disputed outcome would arise. & 0.569 {\tiny$[0.535, 0.604]$} & $<$0.001 \\
\bfseries Meta-level epistemic challenge & The challenger contests the framing, method, evidence standards, or burden of proof rather than only the first-order claim. & 0.561 {\tiny$[0.527, 0.596]$} & $<$0.001 \\
\bfseries Reframing and normalization & The challenger tries to shift the evaluative lens by redescribing the issue, normalizing it, or offering an alternative way to view it. & 0.561 {\tiny$[0.527, 0.596]$} & $<$0.001 \\
\bfseries Practical consequences and feasibility$^{*}$ & The challenger evaluates the view in terms of real-world effects, costs, risks, tradeoffs, or implementability. & 0.558 {\tiny$[0.524, 0.593]$} & $<$0.001 \\
\bfseries Targeted rebuttal & The challenger engages the OP's specific claims, premises, or subpoints directly and often one by one. & 0.554 {\tiny$[0.520, 0.589]$} & $<$0.001 \\
\bfseries Exchange empirical framing$^{*}$ & The exchange is framed in empirical or quantitative terms, emphasizing measurement, data, or evidentiary specificity. & 0.553 {\tiny$[0.519, 0.588]$} & 0.002 \\
\bfseries Alternative proposal and guidance$^{*}$ & The challenger offers concrete substitute policies, remedies, or practical steps instead of only criticizing. & 0.547 {\tiny$[0.513, 0.582]$} & 0.006 \\
\bfseries Rhetorical vividness & The argument uses especially vivid phrasing, storytelling, or memorable imagery to make the point stick. & 0.546 {\tiny$[0.512, 0.581]$} & 0.007 \\
\bfseries Historical framing and legacy & The argument appeals to historical context, precedent, innovation, or long-run cultural legacy. & 0.545 {\tiny$[0.510, 0.580]$} & 0.003 \\
Audience tailoring & The challenger tailors the argument to particular subgroups, audiences, or the OP's specific situation. & 0.534 {\tiny$[0.499, 0.568]$} & 0.0502 \\
Concessionary rebuttal & The challenger softens resistance by explicitly granting part of the OP's point before disagreeing. & 0.532 {\tiny$[0.498, 0.567]$} & 0.070 \\
Consistency challenge$^{*}$ & The reply pressures the OP by showing that their own standard, principle, or rule yields inconsistency or self-undermining conclusions. & 0.524 {\tiny$[0.489, 0.558]$} & 0.195 \\
Legal and institutional framing & The challenger grounds the rebuttal in law, institutions, official roles, or procedural distinctions. & 0.516 {\tiny$[0.481, 0.550]$} & 0.369 \\
Emotional appeal and empathy & The reply persuades through feelings, empathy, reassurance, or existential/personal significance. & 0.516 {\tiny$[0.481, 0.550]$} & 0.412 \\
Psychological framing and diagnosis & The challenger interprets views or behavior in terms of motives, traits, psychology, or developmental state. & 0.502 {\tiny$[0.467, 0.537]$} & 0.966 \\
Normative and rights appeal$^{*}$ & The challenger argues from moral principle, fairness, rights, duty, or harm rather than only from descriptive claims. & 0.501 {\tiny$[0.466, 0.535]$} & 1.000 \\
Challenger experiential grounding$^{*}$ & The challenger or exchange relies on firsthand experience, self-disclosure, or identity-based testimony as support. & 0.497 {\tiny$[0.462, 0.532]$} & 0.907 \\
Rhetorical questioning & The challenger presses the point through pointed questions or requests for clarification rather than only direct assertion. & 0.494 {\tiny$[0.459, 0.528]$} & 0.773 \\
Confrontational tone$^{*}$ & The reply relies on combative, loaded, mocking, or emotionally aggressive rhetoric. & 0.485 {\tiny$[0.450, 0.520]$} & 0.427 \\
\bottomrule
\end{tabularx}
% 19 of 29 features have q < 0.05; * = D_r (16 features)
